% ============================================================
%  Capítulo 2 — Fundamentação Teórica e Trabalhos Relacionados
% ============================================================
\chapter{FUNDAMENTAÇÃO TEÓRICA E TRABALHOS RELACIONADOS}
\label{cap:fundamentacao}

\section{Browser fingerprinting}
\label{sec:fingerprinting}

Browser fingerprinting é a prática de coletar atributos do navegador e do
dispositivo do usuário para construir um identificador persistente sem o uso de
mecanismos de armazenamento local como cookies \cite{eckersley2010}. Ao
contrário dos cookies, um fingerprint não deixa rastros no sistema
de arquivos do usuário, não pode ser apagado com a limpeza do histórico e persiste
mesmo em modo de navegação privada. Estudos de larga escala demonstram que a
combinação de vetores de fingerprinting é suficientemente discriminativa
para identificar unicamente mais de 99\% dos visitantes de um mesmo site
\cite{laperdrix2016}.

A coleta de fingerprints ocorre inteiramente em JavaScript, executado
dentro do processo do navegador. As principais superfícies de ataque são descritas
a seguir.

\textbf{Canvas.} A API Canvas 2D permite que scripts desenhem formas, texto e
gradientes em uma área de renderização virtual. Pequenas diferenças na pilha
gráfica do sistema operacional, motor de renderização, antialiasing, driver de
GPU, fonte instalada, produzem pixels ligeiramente distintos para o mesmo
comando de desenho. A extração do conteúdo via \texttt{canvas.toDataURL()} ou
\texttt{canvas.getImageData()} fornece uma sequência de bytes que varia entre
dispositivos e navegadores \cite{mowery2012}. A mesma superfície é explorada pela
font fingerprinting: chamadas a \texttt{canvas.measureText()} com
diferentes famílias tipográficas revelam quais fontes estão instaladas no sistema.

\textbf{Áudio.} A Web Audio API expõe um grafo de processamento de sinais.
Scripts de fingerprinting instanciam um \texttt{OscillatorNode}, conectam
um \texttt{DynamicsCompressorNode} e renderizam o resultado em modo offline.
As características do sinal resultante, determinadas pela implementação de
aritmética de ponto flutuante do motor de áudio e pelas configurações do driver de
áudio, diferem entre plataformas e arquiteturas de CPU \cite{englehardt2016}.

\textbf{WebGL.} A API WebGL expõe parâmetros da GPU via
\texttt{gl.getParameter()} e \texttt{gl.getSupportedExtensions()}. A combinação de
modelo de GPU, versão de driver e extensões suportadas forma um vetor altamente
discriminativo, particularmente em ambientes corporativos onde a diversidade de
hardware é limitada.

\textbf{WebRTC.} O protocolo WebRTC, destinado à comunicação em tempo real entre
navegadores, exige a negociação de candidatos ICE (Interactive Connectivity
Establishment), que incluem endereços IP locais da máquina. Scripts maliciosos
podem utilizar \texttt{RTCPeerConnection.createOffer()} e
\texttt{setLocalDescription()} para recuperar o endereço IP local do usuário mesmo
quando este utiliza uma VPN, tornando o WebRTC um vetor de identificação de rede
\cite{reiter2017}.

\textbf{Tela e propriedades de navegador.} As dimensões e profundidade de cor da
tela (\texttt{screen.width}, \texttt{screen.height}, \texttt{screen.colorDepth},
\texttt{screen.pixelDepth}) são facilmente acessíveis sem qualquer API
especializada. O objeto \texttt{navigator} expõe propriedades de hardware e
configuração do sistema: \texttt{hardwareConcurrency} (número de núcleos lógicos),
\texttt{deviceMemory} (memória RAM em GB, discretizada), \texttt{languages}
(idiomas configurados), \texttt{platform} e \texttt{maxTouchPoints}. Isoladas,
essas propriedades têm baixa entropia; combinadas com outros vetores, contribuem
significativamente para a unicidade do fingerprint \cite{laperdrix2016}.

\textbf{Sensores e APIs emergentes.} Em dispositivos móveis, APIs como
\texttt{DeviceMotionEvent}, \texttt{DeviceOrientationEvent} e \texttt{Gyroscope}
expõem dados de sensores físicos que podem ser utilizados para
fingerprinting mesmo sem permissão explícita do usuário \cite{das2018}.
Embora menos relevantes para a configuração de desktop estudada neste
trabalho, esses vetores exemplificam a tendência de expansão contínua da superfície
de fingerprinting à medida que novas APIs são adicionadas aos padrões web.

A presença de bibliotecas de fingerprinting de código aberto, FingerprintJS, CreepJS e variantes, nos mil sites mais populares da web indica
que o rastreamento por fingerprint é amplamente adotado na prática. Essas
bibliotecas empacotam os vetores descritos acima em funções JavaScript que executam
silenciosamente durante o carregamento da página, sem interação ou consentimento do
usuário.

%, , , , , , , , , , , , , , , , , , , --
\section{Arquitetura dos motores JavaScript}
\label{sec:motores}

Motores JavaScript modernos compilam código-fonte para uma representação em
bytecode que é interpretada por uma máquina virtual, com caminhos de compilação JIT
(just-in-time) para funções de alta frequência. As duas implementações mais
utilizadas em navegadores de desktop, V8 (Google) e SpiderMonkey (Mozilla)
compartilham esse modelo geral, mas diferem substancialmente em sua arquitetura
interna, o que resulta em sequências de bytecodes distintas para o mesmo código
JavaScript.

\subsection{V8 e o interpretador Ignition}
\label{sub:v8}

O V8 é o motor JavaScript do Chromium e do Node.js. Seu pipeline de compilação
opera em três estágios principais. O analisador sintático converte o código-fonte
em uma árvore sintática abstrata (AST). O gerador de bytecodes Ignition percorre a
AST e emite sequências de instrução para uma máquina de registradores virtual com
acumulador implícito. O compilador otimizador TurboFan monitora, em tempo de
execução, as funções mais frequentemente executadas e recompila seletivamente as
hot functions para código de máquina nativo.

O conjunto de instruções do Ignition inclui operações de carga e armazenamento em
registradores (\texttt{Lda*}, \texttt{Sta*}), operações aritméticas e de
comparação, chamadas de função (\texttt{Call*}, \texttt{Invoke*}), acesso a
propriedades (\texttt{GetNamedProperty}, \texttt{GetKeyedProperty}), controle de
fluxo (\texttt{Jump*}) e operações de closure e contexto (\texttt{CreateClosure},
\texttt{LdaContextSlot}). Cada instrução é acompanhada de metadados de
feedback utilizados pelo TurboFan para inferência de tipos.

A instrumentação deste trabalho é inserida na função \texttt{DoFinalizeJobImpl} do
arquivo \path{interpreter.cc}, que é invocada ao final da compilação de cada
função. Nesse ponto, a sequência completa de bytecodes da função está disponível e
pode ser emitida para o stderr junto com metadados de identificação: URL do
script, posições de início e fim no código-fonte, e nome da função.

\subsection{SpiderMonkey}
\label{sub:spidermonkey}

O SpiderMonkey é o motor JavaScript do Firefox. Seu pipeline segue uma arquitetura
semelhante, análise sintática, geração de bytecodes, interpretação e compilação
JIT, mas difere do V8 em aspectos relevantes para este trabalho.

O formato de bytecodes do SpiderMonkey é baseado em pilha (stack-based), ao
contrário do formato orientado a registradores do Ignition. Isso implica que
operações que o V8 codifica em uma única instrução de três endereços frequentemente
correspondem a sequências de instruções de empilhamento e desempilhamento no
SpiderMonkey. Por exemplo, a leitura de uma propriedade nomeada em V8 emite
\texttt{GetNamedProperty} com dois operandos explícitos, enquanto o SpiderMonkey
emite \texttt{GetProp} precedido de uma instrução de carga do objeto receptor.

Adicionalmente, o SpiderMonkey utiliza estratégias distintas de alocação de
registradores, inlining de funções e representação de closures. A consequência
prática é que mesmo após a abstração dos opcodes para uma representação
intermediária compartilhada (descrita no Capítulo~\ref{cap:ir}), as distribuições
de n-gramas de IR geradas por ambos os motores para o mesmo código JavaScript
diferem estruturalmente, fato que motiva a necessidade de treinamento conjunto
documentada no Capítulo~\ref{cap:experimentos}.

A instrumentação do SpiderMonkey neste trabalho emite blocos de múltiplas linhas
para o stderr com um prefixo identificável (\texttt{[SpiderMonkey Log]}),
capturados via log do geckodriver. Para funções anônimas, que representam a
maioria das funções de fingerprinting em bundles minificados, um identificador único é injetado como \texttt{displayName} da função,
propagando-se para \texttt{Error.stack} e permitindo o join com as chamadas
de API sem perda de dados.

%, , , , , , , , , , , , , , , , , , , --
\section{ByteDefender: abordagem, resultados e limitações}
\label{sec:bytedefender}

ByteDefender \cite{nikkhahbahrami2025} é a referência mais diretamente relacionada
ao presente trabalho. Ele demonstrou pela primeira vez que sequências de bytecodes do
V8 são discriminativas o suficiente para detectar funções de fingerprinting
a nível de função com alta precisão, sem acesso ao código-fonte original.

\textbf{Abordagem.} O ByteDefender instrumenta o motor V8 para capturar as
sequências de bytecodes produzidas por cada função compilada. As sequências são
representadas como bag-of-ngrams de opcodes e utilizadas como entrada para
um classificador baseado em Transformer. O treinamento é supervisionado com rótulos
derivados de heurísticas de correspondência entre chamadas de APIs de
fingerprinting (obtidas por instrumentação de APIs em nível de JavaScript) e
as funções que as invocaram.

\textbf{Resultados.} O ByteDefender reportou F1 $\approx 0{,}95$ na detecção
binária de funções de fingerprinting sobre um corpus de sites da web
visitados com o Chromium.

\textbf{Limitações.} O ByteDefender opera exclusivamente sobre bytecodes do V8.
Modelos treinados sobre bytecodes brutos do V8 não são diretamente aplicáveis ao
SpiderMonkey: os conjuntos de opcodes são disjuntos entre os dois motores, e não
existe mapeamento automático entre eles. A adição de suporte a um novo motor requer
a coleta de um novo corpus de bytecodes, a adaptação do pipeline de dados e o
retreinamento completo do modelo, uma barreira de engenharia não trivial.

Este trabalho estende o ByteDefender em três dimensões: (1) adiciona instrumentação
do SpiderMonkey com rastreamento de funções anônimas; (2) projeta uma representação
intermediária que unifica os vocabulários de bytecodes de ambos os motores; e (3)
demonstra que um único modelo treinado conjuntamente sobre essa representação
generaliza para ambos os motores sem retreinamento, ao custo de aproximadamente 5
pontos de F1 em relação à abordagem de motor único (F1~=~0,904 vs.\
$\approx$0,950), custo que inclui também a inclusão de funções anônimas, um
espaço de rótulos mais amplo e a ausência de metadados de nome de função.

%, , , , , , , , , , , , , , , , , , , --
\section{NLP aplicado a código: TF-IDF, embeddings e representações intermediárias}
\label{sec:nlp}

O tratamento de sequências de bytecodes como texto é a ponte conceitual entre
compilação e aprendizado de máquina. Esta seção descreve as técnicas de
representação e classificação utilizadas no trabalho.

\subsection{TF-IDF}
\label{sub:tfidf}

TF-IDF (term frequency--inverse document frequency) é um esquema de
ponderação clássico para representação de documentos em espaços vetoriais
\cite{manning2008}. Dado um corpus de documentos, o peso de um termo $t$ em um
documento $d$ é o produto da frequência do termo no documento
($\text{tf}(t,d)$) pelo logaritmo inverso da proporção de documentos que contêm o
termo ($\text{idf}(t)$). Termos ubíquos recebem peso baixo; termos discriminativos
recebem peso alto.

Neste trabalho, cada função JavaScript é tratada como um documento e cada token de
RI (ou n-grama de tokens) como um termo. A representação TF-IDF de unigrams produz
vetores de 136 dimensões; a de bigrams, 4.723; a de trigrams, 50.631. A
escalabilidade do TF-IDF a conjuntos de features muito maiores que o vocabulário de
entrada, por conta da esparsidade das representações em bag-of-ngrams, é uma de suas vantagens práticas sobre embeddings densos no cenário de vocabulário
fechado de 145 tokens estudado neste trabalho.

\subsection{Word2Vec e FastText}
\label{sub:embeddings}

Word2Vec \cite{mikolov2013} introduziu o aprendizado de representações distribuídas
de palavras por predição de contexto, utilizando as arquiteturas CBOW
(Continuous Bag of Words) e Skip-gram. Dada uma sequência de tokens,
o modelo Skip-gram aprende vetores de forma que tokens com contextos
similares sejam próximos no espaço vetorial.

FastText \cite{bojanowski2017} estende o Word2Vec ao decompor cada palavra em
n-gramas de caracteres, somando seus vetores para compor a representação final.
Isso permite representar palavras fora do vocabulário, uma vantagem em corpora
de código com identificadores raros. Para o vocabulário fechado de 145 tokens deste
trabalho, FastText é utilizado em modo Skip-gram com subsampling de
frequência desativado ($t = 0{,}0$), pois qualquer subsampling eliminaria
completamente o sinal de treinamento com um vocabulário de apenas 145 tokens. Os
vetores resultantes são avaliados como features de classificadores sequenciais no
Capítulo~\ref{cap:experimentos}.

\subsection{Random Forest e LightGBM}
\label{sub:rf-lgbm}

Random Forest \cite{breiman2001} é um método de ensemble que treina múltiplas
árvores de decisão sobre subamostras aleatórias do conjunto de treinamento e de
features, agregando as predições por votação ou média de probabilidades. A
estratégia \texttt{balanced\_subsample} reamostra cada árvore para balancear as
classes, o que confere robustez natural a classes raras, propriedade relevante
para o cenário multilabel estudado no Capítulo~\ref{cap:experimentos}, onde
categorias como WebRTC e Fonte têm apenas dezenas de exemplos positivos.

LightGBM \cite{ke2017} é uma implementação de Gradient Boosted Decision
Trees (GBDT) baseada em histogram-based splitting e leaf-wise tree
growth. Sua eficiência computacional com features esparsas de alta dimensionalidade
como os vetores TF-IDF de bigrams (4.723 features) e trigrams (50.631 features)
utilizados neste trabalho, torna-o adequado para conjuntos de dados de escala
industrial. LightGBM inclui o parâmetro \texttt{is\_unbalance}, que ajusta
automaticamente os pesos de classe na função de perda para lidar com
desbalanceamento.

\subsection{SVM com SGD e calibração de probabilidades}
\label{sub:svm-fund}

Máquinas de Vetores de Suporte (SVM) buscam o hiperplano que maximiza a margem
de separação entre classes no espaço de features. Para representações esparsas
de alta dimensão como os vetores TF-IDF de bigrams utilizados neste trabalho,
o treinamento via gradiente estocástico (\textit{Stochastic Gradient Descent},
SGD) com perda de dobradiça (\textit{hinge loss}) é computacionalmente
eficiente, com complexidade linear no número de amostras \cite{manning2008}.
O SVM com SGD produz pontuações de margem, não probabilidades; a conversão
para probabilidades necessária ao ajuste de limiar é feita por calibração de
Platt (\texttt{CalibratedClassifierCV}), que aplica regressão logística sobre
as pontuações de margem em validação cruzada interna.

\subsection{BiLSTM}
\label{sub:bilstm-fund}

Redes de Memória de Longo Prazo Bidirecional (\textit{Bidirectional Long
Short-Term Memory}, BiLSTM) processam sequências em ambas as direções, direta
e reversa, produzindo para cada posição uma representação que capta contexto
anterior e posterior. Uma BiLSTM empilha dois RNNs: uma camada que percorre a
sequência da esquerda para a direita e uma camada reversa; as saídas são
concatenadas posição a posição. A representação da sequência completa é obtida
por \textit{pooling} sobre os vetores de saída; neste trabalho utiliza-se
\textit{mean+max pooling}, que concatena a média e o máximo elemento a elemento
ao longo da dimensão temporal, capturando tanto a distribuição central quanto
os picos de ativação mais informativos da sequência de tokens de RI.

\subsection{Transformers}
\label{sub:transformers}

A arquitetura Transformer \cite{vaswani2017} introduziu o mecanismo de
auto-atenção (self-attention), que pondera cada elemento de uma sequência em
função de todos os demais, capturando dependências de longo alcance sem recorrência.
O bloco básico do Transformer empilha camadas de atenção multi-cabeça seguidas de
redes feedforward posição a posição, com normalização de camada e conexões
residuais. Desde sua proposição, a arquitetura se tornou o padrão dominante em
tarefas de processamento de linguagem natural e, progressivamente, em classificação
de código-fonte.

Neste trabalho, um classificador Transformer de pequeno porte é avaliado em duas
configurações: com embeddings FastText fixos e com embeddings iniciados aleatoriamente
e treinados conjuntamente com o classificador. Os resultados e a comparação com os
classificadores baseados em TF-IDF estão no Capítulo~\ref{cap:experimentos}.

\subsection{Representações intermediárias em compiladores}
\label{sub:ri-compiladores}

O conceito de representação intermediária (RI) é central na teoria de compiladores:
uma RI é uma forma de programa independente da linguagem de entrada e da
arquitetura-alvo, sobre a qual transformações e análises podem ser aplicadas de
forma genérica \cite{aho2006}. O exemplo mais influente na prática é o LLVM IR
\cite{lattner2004}, uma representação em forma SSA (static single
assignment) que permite que front-ends de múltiplas linguagens e
back-ends de múltiplas arquiteturas compartilhem um mesmo conjunto de
otimizações.

Este trabalho adota um princípio análogo ao nível dos bytecodes: em vez de operar
sobre os conjuntos de opcodes proprietários do V8 e do SpiderMonkey, projeta-se uma
RI específica para fingerprinting que abstrai diferenças de implementação
entre os dois motores, preservando as distinções semânticas relevantes para a
tarefa de classificação. A RI não é gerada automaticamente por um compilador, mas
construída empiricamente por meio de micro-scripts que exercitam construções
JavaScript relevantes e do alinhamento manual e estatístico das sequências de
bytecodes produzidas por ambos os motores (detalhado no Capítulo~\ref{cap:ir}).

%, , , , , , , , , , , , , , , , , , , --
\section{Trabalhos relacionados}
\label{sec:relacionados}

Esta seção posiciona o presente trabalho no contexto da literatura existente sobre
detecção de fingerprinting e análise de JavaScript.

\textbf{Detecção comportamental de fingerprinting.} \citeonline{iqbal2021}
propõem o FP-Inspector, um sistema que combina análise estática de JavaScript com
aprendizado de máquina para detectar scripts de fingerprinting a nível de
script (não de função). O trabalho usa features extraídas da AST e do fluxo de
controle, atingindo alta precisão na classificação de scripts completos, mas sem
granularidade funcional. Nossa abordagem opera no nível de função individual e usa
bytecodes em vez de código-fonte, tornando-a robusta à ofuscação.

\textbf{Detecção de JavaScript de rastreamento por bytecodes.} O trabalho mais
próximo do nosso em metodologia é apresentado por \citeonline{ghasemisharif2023} em
``Read Between the Lines: Detecting Tracking JavaScript with Bytecode
Classification''. Os autores utilizam bytecodes do V8 para classificar scripts de
rastreamento (tracking), mostrando que a representação em bytecodes é
discriminativa mesmo quando o código-fonte é ofuscado. Diferimos desse trabalho em
dois aspectos principais: operamos no nível de função (granularidade mais fina), e
introduzimos uma RI que permite a generalização cross-engine.

\textbf{Fingerprinting de canvas.} \citeonline{mowery2012} em
``Pixel Perfect: Fingerprinting Canvas in HTML5'' descrevem detalhadamente
o vetor Canvas e demonstram sua eficácia para identificação única de navegadores.
Este trabalho é diretamente relevante para a categoria Canvas do nosso classificador
multilabel.

\textbf{WebRTC como vetor de identificação.} \citeonline{reiter2017} documentam em
``WebRTC: Your Privacy is at Risk'' o mecanismo pelo qual candidatos ICE
expõem IPs locais, constituindo o vetor WebRTC da nossa taxonomia de
fingerprinting.

\textbf{Sensores em dispositivos móveis.} \citeonline{das2018} em ``The
Web's Sixth Sense: A Study of Scripts Accessing Smartphone Sensors'' estudam o
acesso não autorizado a sensores físicos por scripts web. O trabalho demonstra que
APIs de sensor são exploradas para fingerprinting em dispositivos móveis,
expandindo a superfície de rastreamento além dos vetores clássicos.

\textbf{ByteDefender.} \citeonline{nikkhahbahrami2025} em ``Byte by Byte:
Unmasking Browser Fingerprinting at the Function Level Using V8 Bytecode
Transformers'' são o trabalho de referência direta para o presente trabalho
(detalhado na Seção~\ref{sec:bytedefender}).

A Tabela~\ref{tab:relacionados} compara as abordagens mais relevantes.

\begin{table}[htb]
\centering
\caption{Posicionamento comparativo dos trabalhos relacionados}
\label{tab:relacionados}
\begin{tabularx}{\linewidth}{>{\raggedright\arraybackslash}X l l >{\raggedright\arraybackslash}X r}
\toprule
\textbf{Trabalho} & \textbf{Nível} & \textbf{Motor(es)} & \textbf{Features} & \textbf{F1 (binário)} \\
\midrule
FP-Inspector \cite{iqbal2021}            & Script  & N/A (AST) & AST + fluxo de controle  & $0{,}935$ \\
Read Between the Lines \cite{ghasemisharif2023} & Script & V8 & Bytecodes brutos & $0{,}944$--$0{,}956$$\ddagger$ \\
ByteDefender \cite{nikkhahbahrami2025}   & Função  & V8        & Bytecodes brutos V8      & $\approx 0{,}950$ \\
\textbf{Este trabalho}                   & \textbf{Função} & \textbf{V8 + SM} & \textbf{RI cross-engine (145 tokens, 136 ativos)} & \textbf{0,853 / 0,904$\dagger$} \\
\bottomrule
\end{tabularx}
{\footnotesize $\dagger$ 0,853: melhor modelo de bytecodes puros (LightGBM + bigramas + densidade de script).
0,904: modelo implantável com inferência iterativa, sem rótulos externos.
$\ddagger$ Três variantes reportadas: RF + Word2Vec (0,944), DPCNN + FastText (0,952), FastText (0,956).}

\fonte{Elaborada pelo autor.}
\end{table}

A principal distinção deste trabalho em relação ao estado da arte é a operação
cross-engine: nenhum trabalho anterior propôs uma representação intermediária que
unifique os bytecodes de múltiplos motores JavaScript para fins de classificação de
fingerprinting, nem avaliou quantitativamente o custo de F1 dessa
generalização.
